\section{Problem 654: a precise defect parameter for the remaining lower bound}
\subsection*{1) OBJECTIVE AND SOURCE REVIEW}
Attempt dated 2026-09-23. I read all of \texttt{654.tex}. Its two-axis construction correctly addresses the already disproved near-$n$ assertion under no four concyclic points, but does not settle the requested improvement above $(n-1)/3$, including the no-three-collinear version. The primary Aletheia paper likewise presents this as a partial result. I investigate the remaining general-position lower bound, rather than promoting the known counterexample to a full solution.
\subsection*{2) WORK}
Let $X$ have $n$ distinct points, no three collinear and no four concyclic. For each centre $x\in X$, let $s_i(x)$ count its distance classes of size $i$, for $i=1,2,3$. Larger classes are excluded by the no-four-concyclic hypothesis. Thus
\[
n-1=s_1(x)+2s_2(x)+3s_3(x),\qquad
R_X(x)=s_1(x)+s_2(x)+s_3(x).
\]
For an unordered pair $\{y,z\}$, let $c(y,z)$ be the number of points $x\in X$ equidistant from $y,z$. Such points lie on the perpendicular bisector, so $0\le c(y,z)\le2$. Define its total deficit
\[
D(X)=\sum_{\{y,z\}\subset X}(2-c(y,z))\ge0.
\]
Double-counting choices of an apex and an equal-distance pair yields
\[
\sum_{\{y,z\}}c(y,z)=\sum_x(s_2(x)+3s_3(x)).
\]
Subtract from $n(n-1)=\sum_x(s_1+2s_2+3s_3)$ to obtain the exact identity
\[
D(X)=\sum_x(s_1(x)+s_2(x)).
\]
Since $3R_X(x)-(n-1)=2s_1(x)+s_2(x)$, averaging now gives
\[
\max_xR_X(x)\ge\frac{n-1}{3}+\frac{D(X)}{3n}.
\]
In particular, it would suffice to prove a uniform $D(X)\ge\delta n^2$ for some $\delta>0$. More concretely, if at least $\delta n^2$ unordered pairs have at most one equidistant apex in $X$, that same improved lower bound follows.

There is also a direct obstruction to transferring the supplied construction into general position while preserving all its distance coincidences. If three distinct observer points remain equidistant from the same pair of opposite-axis points, those three observers must remain on their common perpendicular bisector. They are therefore collinear. A perturbation satisfying no three collinear must destroy some of precisely the equalities used by the original construction; generic small displacement does not preserve its distance count bound.
\subsection*{3) ADVERSARIAL VERIFICATION}
The formula for $D$ is an exact double count, not evidence that $D$ is quadratic. The elementary restrictions only guarantee nonnegativity, and proving the required quantitative deficit is the unresolved geometric step. The two-axis example contains arbitrarily large collinear subsets and cannot be used in this argument's general-position class.
\subsection*{4) FINAL}
\textbf{UNRESOLVED.} A uniform quadratic deficit of saturated perpendicular bisectors would imply the desired positive linear improvement. No such deficit theorem or matching general-position construction is proved here.
\subsection*{5) SOURCES CHECKED}
The complete supplied version; T. Feng et al., \emph{Semi-Autonomous Mathematics Discovery with Gemini: A Case Study on the Erd\H{o}s Problems}, \texttt{https://arxiv.org/html/2601.22401v3}, checked 2026-09-23 for the attribution and partial scope. The defect identity and perturbation obstruction are proved above.
\subsection*{6) COMPLETION ESTIMATE}
Subjective progress toward the remaining linear-improvement question.
\textbf{COMPLETION: 15\%}
